\chapter{Проективный ёж, точечный дефект и необходимость спинорного подъёма}

Предыдущий гейт нашёл условную нейтринную линию после совместного действия
голономии и Хиггса. Но одномерный кинк в трёх пространственных измерениях
становится стенкой, а не частицей. Поэтому перед построением общего
операторнозначного функционала требуется проверить коразмерность дефекта.

Настоящая глава сравнивает скалярную стенку с проектно-семейным точечным
дефектом, уже подсказанным ранг-один геометрией проекта.

\section{Коразмерность скалярного кинка}

Вакуумное множество вещественного поля с двойной ямой состоит из двух
компонент. Нетривиальный элемент $\pi_0$ защищает дефект коразмерности
один. В одном пространственном измерении это локализованный кинк, но в
$3+1$ измерениях тот же профиль описывает двумерную доменную стенку.

Умножение массы на семейный проектор и хиггсовское направление не меняет
эту геометрию: выражение
\begin{equation}
 q(z)P_0\otimes P_\nu(H)
\end{equation}
локализует фермион на стенке, а не в точке. Поэтому прежний
proof-of-concept подтверждает механизм дефекта, но не форму элементарной
частицы в трёхмерном пространстве.

Минимальное хиггсовское вакуумное многообразие Стандартной модели также не
даёт нужной точечной топологии: оно гомотопически эквивалентно $S^3$ и
имеет тривиальную $\pi_2$; см. \path{arXiv:1610.05623}. Хиггс выбирает
слабое направление уже существующего дефекта, но сам по себе не создаёт
топологически устойчивый монопольный класс минимальной модели.

\section{Проекторное вакуумное многообразие}

Ранг-один семейный проектор можно записать как
\begin{equation}
 P(\mathbf n)=\mathbf n\mathbf n^T,
 \qquad \mathbf n\in S^2,
 \qquad \mathbf n\sim-\mathbf n.
 \label{eq:v5-projective-order-parameter}
\end{equation}
Следовательно, пространство проекторных вакуумов равно
\begin{equation}
 \mathcal V_P\simeq\mathbb{RP}^2,
 \qquad
 \pi_2(\mathbb{RP}^2)\simeq\mathbb Z.
 \label{eq:v5-rp2-pi2}
\end{equation}
Вокруг точки $\mathbf x=0$ естественный представитель задаётся
\begin{equation}
 \mathbf n(\mathbf x)=\frac{\mathbf x}{|\mathbf x|},
 \qquad
 P(\mathbf x)=\frac{\mathbf x\mathbf x^T}{|\mathbf x|^2}.
 \label{eq:v5-projective-hedgehog}
\end{equation}
На окружающей сфере его ориентированный подъём имеет степень
\begin{equation}
 \frac1{4\pi}\int_{S^2}
 \mathbf n\cdot
 (\partial_\theta\mathbf n\times\partial_\phi\mathbf n)
 \,d\theta d\phi=1.
 \label{eq:v5-hedgehog-degree}
\end{equation}
Таким образом, проектный порядок способен нести точечный топологический
дефект. Это более подходящая геометрия частицы, чем скалярная стенка.

\section{Энергетическая граница глобального ежа}

Для одного сигма-модельного градиентного члена
\begin{equation}
 E_P=\frac12\int d^3x\,\Tr(\partial_iP\,\partial_iP)
 \label{eq:v5-projector-gradient-energy}
\end{equation}
у ежа выполнено
\begin{equation}
 \Tr(\partial_iP\,\partial_iP)=\frac4{r^2}.
\end{equation}
Энергия каждой радиальной оболочки равна $8\pi,dr$, поэтому полная энергия
растёт линейно с радиусом системы. Это стандартная граница глобальных
монополей: калибровочное поле может компенсировать дальний угловой
градиент, тогда как чистый глобальный ёж имеет расходящуюся энергию; см.
\path{arXiv:gr-qc/0307074}.

Проект действительно содержит семейную связность, но пока она является
внутренней операторной связностью, а не выведенным пространственным
$SO(3)$-полем с кинетическим членом и конечной монопольной энергией.
Отождествлять их без нового функционала нельзя.

\section{Проекторная масса и нулевой индекс}

Самый прямой оператор, построенный только из $P$, имеет вид
\begin{equation}
 Q_P=3P-I_3.
 \label{eq:v5-projector-mass-matrix}
\end{equation}
Его спектр равен $(2,-1,-1)$, поэтому он невырожден на бесконечности и
может служить матричной массой дираковского оператора.

Однако число защищённых нулевых мод определяется не только степенью
поля порядка, но и топологией собственных расслоений массовой матрицы.
Индекс Каллиаса выражается через классы Черна соответствующих
собственных расслоений; см. \path{arXiv:hep-th/0011081} и
\path{arXiv:hep-th/0311215}.

Положительное собственное пространство $Q_P$ есть комплексification
вещественной линии, натянутой на глобально определённый вектор
$\mathbf n(\theta,\phi)$. Его проектор равен $P$ и имеет
\begin{equation}
 c_1(P)=\frac1{2\pi i}\int_{S^2}
 \Tr\!\left(P[\partial_\theta P,\partial_\phi P]\right)
 d\theta d\phi=0.
 \label{eq:v5-vector-projector-chern-zero}
\end{equation}
Машинный аудит получает тождественно нулевую плотность. Поэтому
проективный ёж топологически нетривиален как карта в $\mathbb{RP}^2$, но
его прямое векторное массовое представление не даёт ненулевого комплексного
индекса фермиона.

\begin{nogocriterion}[Недостаточность одного проекторного ежа]
Нетривиальный заряд $\pi_2(\mathbb{RP}^2)$ не гарантирует фермионную
нулевую моду для произвольного представления. Масса $3P-I_3$ имеет
тривиальное положительное комплексное собственное расслоение и нулевой
первый класс Черна. Следовательно, ранг-один проектор сам по себе не
замыкает точечную частицу с единичным индексом.
\end{nogocriterion}

\section{Спинорный подъём и единичный индекс}

Если ориентированный вектор $\mathbf n$ доступен, можно построить
двухкомпонентную массовую матрицу
\begin{equation}
 Q_{1/2}(\mathbf n)=\mathbf n\cdot\boldsymbol\sigma.
 \label{eq:v5-spinor-hedgehog-mass}
\end{equation}
Проектор её положительной линии равен
\begin{equation}
 P_+(\mathbf n)=\frac12(I+\mathbf n\cdot\boldsymbol\sigma).
\end{equation}
Прямое вычисление даёт
\begin{equation}
 \frac1{2\pi i}\int_{S^2}
 \Tr\!\left(P_+[\partial_\theta P_+,\partial_\phi P_+]\right)
 d\theta d\phi=1.
 \label{eq:v5-spinor-projector-chern-one}
\end{equation}
Именно такой собственный пучок способен дать единичный индекс и
локализованную нулевую моду точечного дефекта. Общая связь топологического
класса пространственно меняющегося гамильтониана с защищёнными модами
точечных и линейных дефектов систематизирована в
\path{arXiv:1006.0690}.

Но $Q_{1/2}$ не является функцией одного $P$. Замена
$\mathbf n\mapsto-\mathbf n$ сохраняет $P$, но меняет знак
$Q_{1/2}$. Поэтому требуется ориентированный двойной подъём
\begin{equation}
 \mathbb{RP}^2\longleftarrow S^2
\end{equation}
или эквивалентная корневая/спинорная линия, хранящая потерянный знак.

\section{Связь с четырёхтактной фазой}

Проект уже обнаружил корневую фазу $h=\pm i$ с
\begin{equation}
 h^2=-1,
 \qquad h^4=1.
\end{equation}
Такая линия является естественным кандидатом на спинорное дополнение
проекторного направления: после проективного обхода знак меняется, а после
двойного обхода полное состояние возвращается. Однако прежний exhaustive
гейт доказал, что корневая линия не назначена обязательным образом пакету
$H_{15}$.

Следовательно, новая конструктивная цель стала точной. Нужен не ещё один
скалярный потенциал, а функториальный подъём
\begin{equation}
 (P,\text{корневая линия})
 \longmapsto
 \mathbf n\cdot\boldsymbol\sigma
 \label{eq:v5-projector-spinor-lift-target}
\end{equation}
совместно с пространственной семейной связностью, делающей энергию ежа
конечной. Без обоих элементов точечная частица остаётся условной.

\begin{theorem}[Точная форма недостающей архитектуры]
Проекторная семейная геометрия имеет правильную коразмерность и допускает
точечный заряд $\pi_2(\mathbb{RP}^2)=\mathbb Z$. Но чистый глобальный ёж
имеет линейно расходящуюся энергию, а векторная масса $3P-I_3$ имеет
нулевой комплексный индекс. Единичный индекс появляется в спинорном
представлении $\mathbf n\cdot\boldsymbol\sigma$, которое требует
ориентированного двойного подъёма, не определяемого одним $P$.
\end{theorem}

\section{Статус}

Поиск единого операторнозначного функционала теперь ограничен двумя
обязательными новыми связями:
\begin{enumerate}
 \item пространственная семейная связность должна компенсировать дальний
 градиент проекторного ежа и дать конечную энергию;
 \item корневая линия порядка четыре должна функториально поднять
 проекторную ось к спинорной массе с ненулевым индексом.
\end{enumerate}
Хиггс-зависимый проектор затем может выбрать слабую нейтринную компоненту,
но не заменяет эти два шага.

\begin{protocol}
\begin{enumerate}
 \item скалярный кинк как точечная частица в $3+1$:
 \textbf{нет, это стенка};
 \item минимальный Хиггс как источник $\pi_2$: \textbf{нет};
 \item проекторное вакуумное пространство: \textbf{$\mathbb{RP}^2$};
 \item точечный проекторный заряд: \textbf{$1$};
 \item энергия глобального ежа: \textbf{линейно расходится};
 \item класс Черна векторного проекторного представления:
 \textbf{$0$};
 \item класс Черна спинорного представления: \textbf{$1$};
 \item спинорный подъём из одного $P$: \textbf{невозможен без знака};
 \item корневая линия как кандидат подъёма: \textbf{допустима, не назначена};
 \item физическое замыкание: \textbf{не достигнуто}.
\end{enumerate}
\end{protocol}

Машинный протокол сохранён в
\path{s2t_v5_projective_hedgehog_point_defect_gate_results.json}.